\section{Aturan Inferensi Kuantor: Instansiasi dan Jeneralisasi}

Sama seperti logika proposisi yang memiliki aturan inferensi (Modus Ponens, Modus Tollens, dan lain-lain), logika predikat memiliki aturan inferensi khusus yang dirancang untuk menangani simbol kuantifikasi. Aturan-aturan ini memungkinkan kita ``membuka'' kuantor (instansiasi) atau ``memasang'' kuantor (jeneralisasi) \cite{rosen2019}.

\subsection{Universal Instantiation (UI)}

\begin{definisi}[Universal Instantiation]
Jika $\forall x\, P(x)$ bernilai benar untuk seluruh elemen domain $\mathcal{U}$, maka untuk setiap konstanta tertentu $c \in \mathcal{U}$, kita berhak menyimpulkan $P(c)$.

Secara formal:
\[ \frac{\forall x\, P(x)}{P(c)} \quad \text{untuk setiap } c \in \mathcal{U} \]
\end{definisi}

UI adalah aturan yang paling sering digunakan dalam penalaran berbasis basis pengetahuan, karena memungkinkan kita menurunkan fakta spesifik dari aturan umum.

\begin{contoh}
Diberikan aturan: $\forall x\, (Manusia(x) \rightarrow Mortal(x))$ dan fakta $Manusia(Socrates)$.

Langkah inferensi:
\begin{enumerate}
    \item $\forall x\, (Manusia(x) \rightarrow Mortal(x))$ \hfill (aturan KB)
    \item $Manusia(Socrates) \rightarrow Mortal(Socrates)$ \hfill (UI, $x = Socrates$)
    \item $Manusia(Socrates)$ \hfill (fakta KB)
    \item $Mortal(Socrates)$ \hfill (Modus Ponens pada 2 dan 3)
\end{enumerate}
\end{contoh}

\subsection{Universal Generalization (UG)}

\begin{definisi}[Universal Generalization]
Jika $P(c)$ dapat dibuktikan benar untuk konstanta sembarang $c$ yang tidak memiliki asumsi khusus (yaitu, bukti berlaku untuk elemen apa pun dari domain), maka kita sah menyimpulkan $\forall x\, P(x)$.

\[ \frac{P(c) \text{ untuk } c \text{ sembarang}}{\forall x\, P(x)} \]
\end{definisi}

\begin{catatan}
Syarat ``$c$ sembarang'' sangat ketat: $c$ tidak boleh memiliki sifat khusus yang membedakannya dari elemen domain lainnya. Jika bukti $P(c)$ menggunakan asumsi tambahan tentang $c$ (misalnya ``$c$ genap''), maka generalisasi ke $\forall x\, P(x)$ tidak valid.
\end{catatan}

\subsection{Existential Instantiation (EI)}

\begin{definisi}[Existential Instantiation]
Jika $\exists x\, P(x)$ bernilai benar, maka terdapat setidaknya satu elemen $c$ dalam domain sehingga $P(c)$ benar. Kita dapat memperkenalkan konstanta baru $c$ (yang belum digunakan sebelumnya) untuk merujuk pada entitas tersebut.

\[ \frac{\exists x\, P(x)}{P(c)} \quad \text{untuk konstanta baru } c \]
\end{definisi}

Berbeda dengan UI, konstanta $c$ dalam EI bersifat ``saksi'' (\textit{witness}) --- ia merujuk pada satu entitas spesifik (meskipun kita mungkin tidak tahu persisnya siapa) yang memenuhi $P$.

\subsection{Existential Generalization (EG)}

\begin{definisi}[Existential Generalization]
Jika $P(c)$ bernilai benar untuk suatu konstanta tertentu $c \in \mathcal{U}$, maka kita sah menyimpulkan $\exists x\, P(x)$.

\[ \frac{P(c)}{\exists x\, P(x)} \]
\end{definisi}

\begin{contoh}
Dari fakta $Manusia(Socrates)$, kita dapat menyimpulkan $\exists x\, Manusia(x)$ (``terdapat sesuatu yang merupakan manusia'') melalui EG.
\end{contoh}

\subsection{Ringkasan Aturan Inferensi Kuantor}

\begin{table}[htbp]
\centering
\renewcommand{\arraystretch}{1.3}
\begin{tabular}{|p{4cm}|p{3.5cm}|p{4.5cm}|}
\hline
\textbf{Aturan} & \textbf{Arah} & \textbf{Syarat Khusus} \\ \hline
Universal Instantiation & $\forall x\, P(x) \Rightarrow P(c)$ & $c$ sembarang elemen domain \\ \hline
Universal Generalization & $P(c) \Rightarrow \forall x\, P(x)$ & $c$ harus sembarang (tanpa asumsi khusus) \\ \hline
Existential Instantiation & $\exists x\, P(x) \Rightarrow P(c)$ & $c$ konstanta baru (saksi) \\ \hline
Existential Generalization & $P(c) \Rightarrow \exists x\, P(x)$ & $c$ elemen spesifik yang memenuhi $P$ \\ \hline
\end{tabular}
\caption{Ringkasan Empat Aturan Inferensi Kuantor}
\end{table}
